\section{Positive Laurent rows and the actual quadratic beds}
\label{sec:positive-row-producer}
The quadratic construction needs positive factors with exponentially
decaying coefficients, uniformly in the finite period. Taking a separate
matrix square root on each torus would not supply such coefficients.
We construct one analytic symbol and then recover its Laurent row.
This is also how the construction preserves its zero-frequency
conventions.

\subsection{A uniform analytic inverse and square root}
Let $I$ be a fixed finite index set of cardinality $r$. A real Laurent
row is a collection $A_{ij}(z)$, $i,j\in I$, $z\in\mathbb Z^4$.
Suppose, for some $\eta,C>0$,
\begin{equation}\label{eq:positive-row-strip}
 \sup_{i,j}\sum_{z\in\mathbb Z^4}
                    e^{\eta|z|_1}|A_{ij}(z)|\le C.
\end{equation}
Its symbol
\[
                  \widehat A(p)=\sum_z A(z)e^{ip\cdot z}
\]
is periodic in each real coordinate and holomorphic whenever
$\max_a|\operatorname{Im}p_a|<\eta$. If the index set is empty
all the statements below are vacuous, so assume $r\ge1$.
On every smaller strip of width $\sigma$, termwise differentiation
is justified by
\[
 |z|_1 e^{\sigma|z|_1}
             \le(\eta-\sigma)^{-1}e^{\eta|z|_1}.
\]
Thus each entry has derivative norm at most
$C/(\eta-\sigma)$, for the maximum norm on momenta. This
also proves local uniform convergence of the differentiated series.
In operator norm the corresponding Lipschitz bound is at most
$rC/(\eta-\sigma)$, since a matrix with entries bounded by $a$
has operator norm at most $ra$.

Suppose in addition that for every real momentum $p$,
\begin{equation}\label{eq:positive-symbol-band}
                 vI\le\widehat A(p)\le MI,
                              \qquad 0<v\le M,
\end{equation}
as Hermitian forms. These bounds are uniform over the whole
real torus. Let $\sigma=\eta/2$, $D=rC/(\eta-\sigma)$,
and choose
\[
       0<\tau\le\min\{\sigma,v/(2D)\}.
\]
For a complex momentum in this strip, put $q=\operatorname{Re}p$.
The segment joining $q$ to $p$ lies in the strip of width $\sigma$,
and therefore
\[
                   \|\widehat A(p)-\widehat A(q)\|\le v/2.
\]
Diagonalizing the Hermitian matrix at $q$ gives
\[
 \|M^{-1}\widehat A(q)-I\|\le1-v/M.
\]
Hence throughout the smaller strip
\begin{equation}\label{eq:positive-binomial-disk}
      \|M^{-1}\widehat A(p)-I\|\le1-v/(2M)=:\rho<1.
\end{equation}
All constants depend only on the original uniform bounds and the
fixed index set.

Put $Z(p)=M^{-1}\widehat A(p)-I$. The normally convergent series
\[
 \widehat F(p)=\sqrt M\sum_{n\ge0}{1/2\choose n}Z(p)^n,
 \qquad
 \widehat B(p)=M^{-1}\sum_{n\ge0}(-Z(p))^n
\]
define holomorphic periodic symbols on this strip. Normal convergence
follows from $\rho<1$ and absolute convergence of the scalar power
series on every smaller closed disk. Cauchy multiplication gives
$\widehat B\widehat A=I$ and $\widehat F^2=\widehat A$.
On the real torus diagonalization gives eigenvalues
$\lambda^{-1}$ and $\sqrt\lambda$, respectively, where
$v\le\lambda\le M$. Consequently
\begin{equation}\label{eq:positive-root-bands}
 M^{-1}I\le\widehat B\le v^{-1}I,
 \qquad \sqrt v I\le\widehat F\le\sqrt M I.
\end{equation}
In particular these are the positive inverse and positive square root,
with no choice of eigenvectors as the momentum varies.

We verify that they are symbols of real exponentially summable rows.
Either symbol is bounded on a still smaller strip, say of width
$\tau_0>0$, by a constant $K$. Its Fourier coefficient is the
integral over $[-\pi,\pi]^4$ against $e^{-ip\cdot z}$, divided
by $(2\pi)^4$. Shift the $a$th contour to imaginary part
$-\tau_1\operatorname{sign}(z_a)$, where $0<\tau_1<\tau_0$.
Cauchy's theorem applies one coordinate at a time; the vertical
sides cancel by periodicity. The coefficient is therefore bounded
by $Ke^{-\tau_1|z|_1}$. For $0<e<\tau_1$,
\[
 \sum_z e^{e|z|_1}|F_{ij}(z)|
 \le K\left(\sum_{n\in\mathbb Z}
                      e^{-(\tau_1-e)|n|}\right)^4<\infty,
\]
and likewise for $B$. The original real row satisfies
$\widehat A(-p)=\overline{\widehat A(p)}$ on real momenta.
Both power series have real coefficients, so their symbols have
the same symmetry. Changing $p$ to $-p$ in the Fourier integral
shows that every recovered coefficient is real.
Absolute Fourier convergence identifies the recovered symbols with
the displayed analytic functions on the real torus.

On a finite torus, periodize these same rows. The discrete Fourier
transform diagonalizes their convolution matrices, with symbols
evaluated at the finite Fourier frequencies. Thus on every period
$n$, simultaneously,
\begin{equation}\label{eq:positive-row-periods}
 B_nA_n=I,\qquad F_n^{\mathsf T}=F_n,\qquad
 F_n^{\mathsf T}F_n=A_n,
 \qquad \sqrt v I\le F_n\le\sqrt M I.
\end{equation}
The same transform transfers the inverse bounds in
\eqref{eq:positive-root-bands}. This includes the zero frequency
and the one-site torus. The construction has used a fixed number
of form labels; no bound involving the growing number of lattice
sites has entered the analytic inverse or root estimate.

\subsection{The amendment of the two-form precision}
Fix a smoothing stage $j\ge4$ and the positive kernel parameter
$k$. For every finite train at this stage, let $H$ and $B$ be the
pulse and dual-pulse rows constructed in
Sections~\ref{sec:ex-lattice-free-kernel-estimates} and
\ref{sec:ex-lattice-free-kernel-components}. Their real symbols are
diagonal, positive and reciprocal. Their uniform strip bounds
therefore also give uniform positive lower bounds: if a diagonal
entry of each row is at most $C$ and their product is one, each
is at least $C^{-1}$. The preceding Fourier argument yields
$HB=BH=I$ on every period, with uniform spectral bounds.

Let $d_1,d_2$ be the exterior differences from one-forms to
two-forms and from two-forms to three-forms; thus $d_2d_1=0$.
For a fixed amendment coefficient $l\ge0$, define
\begin{equation}\label{eq:positive-amendment}
 X=d_2B,\qquad A=B+lX^{\mathsf T}X,\qquad
 G=A^{-1},\qquad R=I+l d_2B d_2^{\mathsf T},\qquad D=R^{-1}.
\end{equation}
These formulas denote compositions of Laurent rows and hence also
their matrices on every finite period. Products and sums preserve
exponential summability, after taking a common strip width.
On real frequencies $A\ge B>0$ and $R\ge I$, because $B>0$.
The analytic inverse construction supplies $G,D$ with uniform
strip and spectral bounds. The constants may depend on $j,k,l$,
but are independent of the train and finite period.

Two identities explain the choice of this amendment. First,
$XHd_1=d_2d_1=0$, so
\[
                       AHd_1=d_1,\qquad Gd_1=Hd_1.
\]
Thus the potential precision is unchanged. Second, expansion of
\eqref{eq:positive-amendment} gives $RX=d_2A$, since $B$ is
symmetric. Multiplication by $D$ on the left and $G$ on the right
therefore gives
\begin{equation}\label{eq:positive-amendment-intertwine}
                              XG=Dd_2.
\end{equation}
Both $G$ and $D$ are symmetric. At zero momentum $d_2=0$ and
$B=I$, so $G(0)=I$; the harmonic convention is retained.

Choose a uniform lower spectral bound $b>0$ for $G$, and allow
any direct precision $0<p\le b/2$. Then
$G-pI\ge(b/2)I$. It has a uniform strip bound, including the
subtracted constant row, so the analytic square-root construction
gives a row $F$ with
\[
               F_n\ge0,\qquad F_n^{\mathsf T}F_n=G_n-pI
\]
for every finite period, with uniform strip bounds. The full root
matrix is
\[
             J_p=\begin{bmatrix}\sqrt p I\\F_n\end{bmatrix},
                         \qquad J_p^{\mathsf T}J_p=G_n.
\]
Together with $H,B,G,A$ this is the actual quadratic bed. Its
inverse, symmetry, positivity and exterior-intertwining properties
have all been verified above. In the top train the pulse is the
identity, so its diagonal bed row is the identity as required.
The admissible upper bound for $p$ and all row bounds are selected
before any lattice volume or train is chosen.
